\subsection{Secants, tangent fibres, and flag lines}
\label{sec:exceptional-reconstruction-fibres}

We describe the scheme-theoretic fibres of the two-component invariant,
including those not covered by the rank-two certificate.  This is a
geometric description of the residual ambiguity of that invariant,
not an assertion that every pair in such a fibre has isomorphic
failure schemes.
Put \(X=\Gr(4,W)\), and let \(X^\times\) be the open where
both projections are nonzero.  Consider the morphism
\[
 \pi:X^\times\longrightarrow\Pj(E_{175})\times\Pj(E_{35}),
       \qquad [w]\longmapsto([Pw],[Qw]).
\]
For \([w]\in X^\times\), set \(A=Pw\), \(B=Qw\), and
\(H_w=\langle\ell_{F,w}:F\in\mathcal I_{\mathrm{Pl},2}\rangle\).
The rank \(r_w=\dim H_w\) is independent of a representative.

\begin{theorem}[Exceptional fibres of the component invariant]
\label{thm:exceptional-component-fibres}
The fibre of \(\pi\) through \([w]\), as a scheme, is the
intersection of \(X\) with \(\Pj(\C A\oplus\C B)\), with
its two endpoints removed.  In coordinates \([a:b]\) on this line
its closed intersection with \(X\) has the following exhaustive forms.
\begin{enumerate}
\item If \(r_w=2\), its saturated ideal is \((b-a)\); it is the
single reduced point \([w]\).
\item If \(r_w=1\), choose a generator
\(\ell=ca+db\) of \(H_w\).  Its saturated ideal is
\(((b-a)(ca+db))\).  It is the length-two divisor
\([1:1]+[d:-c]\), with multiplicity.  The residual point equals
\([w]\) exactly when \(c+d=0\); it is an excluded endpoint exactly
when \(cd=0\).  Otherwise it gives a second, distinct reduced point
of the fibre of \(\pi\).
\item If \(r_w=0\), the whole line lies in \(X\), and the fibre
is \(\mathbf G_m\).  Such a line consists precisely of the
four-planes \(R\) in a fixed flag
\[
 A_3\subset R\subset B_5\subset W,\qquad
          \dim A_3=3,\quad\dim B_5=5.
\]
Equivalently, the pure endpoints are decomposable and the two
four-planes they represent intersect in dimension three.
\end{enumerate}
On the locally closed reduced locus of case (2) with
\(cd(c+d)\ne0\), passage to the residual point is a regular
involution without fixed points.  The differential of \(\pi\) has a
nonzero kernel at \([w_R]\) precisely when
\begin{equation}\label{eq:exceptional-tangent-criterion}
 Qw_R\in \bigwedge^3R\wedge W.
\end{equation}
This is case (3), or the double-point subcase of case (2).
Its kernel is then one-dimensional.  In particular, a distinct
secant pair does not force infinitesimal failure of reconstruction.
\end{theorem}
\begin{proof}
The ambient projection from \(\Pj(E_{175}\oplus E_{35})\) to the
product of projective spaces has fibre
\(\Pj(\C A\oplus\C B)\setminus\{[A],[B]\}\).
Pullback along the closed immersion of \(X\) proves the first
assertion scheme-theoretically.  The Grassmannian ideal is generated
by its quadratic Pluecker relations.  Their restrictions generate
\((b-a)H_w\).  If \(H_w\) is the full two-dimensional space of
linear forms, these forms have no common projective zero and the
ideal sheaf is \((b-a)\).  If \(H_w\) is one-dimensional, the
nonzero quadratic \((b-a)\ell\) generates an already saturated
principal ideal in \(\C[a,b]\), giving its two zeros with their
multiplicities.  If \(H_w=0\), every Pluecker equation vanishes
identically on the line.  This proves the three scheme structures,
including the endpoint and double-point distinctions.

We recall the flag description with its proof.  Choose a point
\(R_0\) on a line in the Grassmannian and a complement to \(R_0\)
in \(W\).  On the Pluecker affine chart of \(R_0\), the line is
parametrized by graphs of \(tM:R_0\to W/R_0\).  Indeed the
one-replacement Pluecker coordinates are its matrix entries and are
linear in \(t\).  The two-replacement coordinates are the two-by-two
minors of \(tM\).  They are also linear functions on the line, so
their quadratic coefficients vanish.  Thus \(\rank M=1\), since
the line is nonconstant.  Its members share \(\ker M\) of dimension
three and lie in \(R_0+\im M\) of dimension five; closing the affine
line gives the asserted flag line.  Conversely the four-planes
between such a flag form a Pluecker line by wedging a fixed basis of
\(A_3\) with a varying vector of \(B_5/A_3\).  Applied to the
endpoints, this also proves the equivalent formulation.

On a chart where one coefficient of a nonzero \(\ell_{F,w}\) is
invertible, its coefficient pair gives the regular map
\[
 [w]\longmapsto[dA-cB].
\]
Rank one makes different choices proportional, so these maps glue.
The hypotheses \(cd(c+d)\ne0\) exclude endpoints and the double
point.  The residual point has the same two projective components;
its line of recombination is the same line, whose intersection with
\(X\) has exactly the two distinct reduced points.  Repeating the
construction returns \([w]\).  This proves regularity and the
involution assertion on the stated reduced stratum.

Finally the ambient projection has a one-dimensional vertical tangent
at \([w]\), represented by \(B\) modulo \(\C w\).  The affine
tangent space to the Pluecker cone at \(w_R\) is
\(\bigwedge^3R\wedge W\), as follows by differentiating a wedge
of four basis vectors of \(R\).  This proves
\eqref{eq:exceptional-tangent-criterion}.  Equivalently set \(a=1\),
\(b=1+t\); each restricted equation has derivative
\(\ell_{F,w}(1,1)\).  All these derivatives vanish exactly when
\(H_w=0\), or when its generator vanishes at \([1:1]\).
The rank-two case cannot have this property.  The vertical tangent
is nonzero because \(A,B\) are independent.  This proves the last
claims.
\end{proof}

\begin{corollary}[Candidate bounds without a genericity certificate]
\label{cor:exceptional-candidate-bound}
On \(X^\times\), away from the flag-line locus, the two-component
invariant has fibres of length at most two.  Failure of the rank-two
certificate can represent either a double point, a genuine second
candidate, or a discarded pure endpoint.  Only the flag-line case
allows a positive-dimensional fibre.  For an abstract failure-scheme
isomorphism these are necessary candidate alternatives after one
common projective change of coordinates.
\end{corollary}
\begin{proof}
Remove the two endpoints from the three closed intersections in the
theorem, and apply Lemma~\ref{lem:common-g-functoriality} for the
last assertion.  No converse from equality of component pairs to
isomorphism of failure schemes is used.
\end{proof}

The theorem describes every fibre configuration on \(X^\times\).
Theorem~\ref{thm:global-schur-recombination} below determines their
incidence with the smooth Jacobian open: only the rank-two case occurs.
No ambient irreducible-component or codimension classification is
implicit in the fibre classification itself.  It separates positive-dimensional flag
ambiguity from tangent degeneracy and from two distinct candidates;
these distinctions are invisible in the single rank-failure condition.

\begin{example}[A nonempty flag-line stratum]
\label{ex:exceptional-flag-line}
In the monomial notation of Appendix~\ref{app:separating-pencils} put
\[
 A=b_1\wedge b_2\wedge b_3\wedge b_5,\qquad
 B=b_1\wedge b_2\wedge b_3\wedge b_4.
\]
These are highest-weight vectors of weights \((4,3,1,0)\) and
\((5,1,1,1)\); hence \(A\in E_{175}\), \(B\in E_{35}\).
Every \(aA+bB\) is the decomposable vector
\(b_1\wedge b_2\wedge b_3\wedge(ab_5+bb_4)\).
Thus the flag-line stratum is nonempty in \(X^\times\), and this
example supplies an entire \(\mathbf G_m\)-fibre.  It is not an
example on \(G_4^\circ\): the first three quadrics have the common
factor \(x_1\), and the fourth quadric has a zero on \(x_1=0\).
\end{example}
